\documentclass{report}
\usepackage{amsmath,amssymb}
\setlength{\parindent}{0mm}
\setlength{\parskip}{1em}
\begin{document}
\begin{center}
\rule{6in}{1pt} \
{\large Hengbin An \\
{\bf Anderson acceleration and application to three temperature energy equations}}

Institute of Applied Physics and Computational Mathematics \\ No 2 \\ Fenghao East Road \\ Haidian District \\ Beijing \\ China 100094
\\
{\tt an\_hengbin@iapcm.ac.cn}\\
Homer F. Walker\\
Xiaowei Jia\end{center}

Anderson acceleration method is an algorithm for accelerating
the convergence of general fixed point iteration or Picard iteration.
The method was first proposed in 1965 to accelerate the convergence
of self-consistent field iteration in electronic structure computation,
and it has being used successfully in that area ever since.
However, until recently, Anderson acceleration method began to attract attention from
the community of numerical analysis and algorithm development and some
other application areas.
Compared with Newton-like method,
one advantage of Anderson acceleration method is that there is no need to
form the Jacobian matrix. Therefore, the method is easy to be implemented.
In this paper, Anderson accelerated Picard method is employed to solve
a kind strong nonlinear radiation diffusion equation, three temperature energy equations.
To improve the robustness of Anderson acceleration method,
two strategies are incorporated into Anderson acceleration method.
One strategy is used so that the iteration satisfying physical constraint.
Another strategy is about monitoring matrix condition number for the
least squares problem in
the implementation of Anderson acceleration so that numerical stability is guaranteed.
Numerical results show that Anderson accelerated Picard method can solve
three temperature energy equation effectively.
Compared to no accelerated Picard method, at least half iteration numbers can be saved by
using Anderson acceleration method.
Comparison for Jacobian-free Newton-Krylov method to Picard method and
Anderson acceleration method is given.


\end{document}
